\documentclass[11pt]{article}
\usepackage[a4paper,margin=27mm]{geometry}
\usepackage{amsmath,amssymb,amsthm,mathtools,hyperref}
\usepackage{xurl}
\usepackage[T1]{fontenc}
\hypersetup{hidelinks,pdfauthor={Rohit Kumar Jha},
  pdftitle={An explicit irrationality-exponent bound for zeta(3) - r zeta(2)}}
\newtheorem{theorem}{Theorem}[section]
\newtheorem{lemma}[theorem]{Lemma}
\newtheorem{proposition}[theorem]{Proposition}
\newcommand{\Q}{\mathbb Q}
\newcommand{\R}{\mathbb R}
\newcommand{\C}{\mathbb C}
\newcommand{\Z}{\mathbb Z}
\newcommand{\ii}{\mathrm i}
\newcommand{\dd}{\,\mathrm d}
\DeclareMathOperator{\sech}{sech}
\DeclareMathOperator{\diag}{diag}
\title{An explicit irrationality-exponent bound for
  $\zeta(3)-r\zeta(2)$\\\large A Lean-verified quantitative extension}
\author{Rohit Kumar Jha}
\date{8 October 2026}
\begin{document}
\maketitle

\begin{abstract}
We give a quantitative extension of Qian Tang's
determinant proof of the irrationality of $\zeta(3)-r\zeta(2)$ for fixed
$r\in\Q$. The new ingredient is a positive-Gram comparison with an
exponential lower bound on its smallest eigenvalue. Combined with the
upstream arithmetic and energy estimates, this proves
$2\leq\mu(\zeta(3)-r\zeta(2))\leq10000$, with an approximation threshold
depending on $r$. A direct determinant comparison avoids any separate
all-minors assumption. The denominator theorem and the exponent corollary
have been verified in Lean and replayed in a fresh Lean-kernel environment;
their axiom closures contain only the standard axioms. The exponent is
not optimized.
\end{abstract}

\section{Statement, inputs, and attribution}

For an irrational real number $x$, its irrationality exponent is
\[
 \mu(x)=\sup\{\nu>0:0<|x-a/q|<q^{-\nu}
                 \text{ for infinitely many reduced }a/q\in\Q\}.
\]
The supremum is interpreted in the extended reals; the estimate proved
below makes it finite for the numbers considered here. Our formal
real-valued supremum is used only after its defining set is proved
nonempty and bounded above. The principal conclusion is
\[
              2\leq\mu(\zeta(3)-r\zeta(2))\leq10000
                  \qquad(r\in\Q\text{ fixed}).
\]

The starting point is Qian Tang's Lean development~\cite{Tang},
\href{https://github.com/dtq1997/zeta2-zeta3-linear-independence/tree/669c92bf7a3728e0de9ad80c373bc29fb5ca3d58}
{\emph{zeta2-zeta3-linear-independence}}, pinned to commit
\begin{center}\texttt{669c92bf7a3728e0de9ad80c373bc29fb5ca3d58}.\end{center}
It builds on Fauzan's determinant method. Its stated endpoint is qualitative
irrationality for every rational $r$, equivalently the linear independence
of $1,\zeta(2),\zeta(3)$ over $\Q$.
The source declares Lean v4.35.0-rc3 and Mathlib revision
\begin{center}\texttt{c55e6e786f49471c72fbddbec5415808896aec1e}.\end{center}

We use its arithmetic-content theorem and prime-edge nonvanishing theorem
as imported mathematical inputs. Its energy proof is strengthened below by
retaining the positive integral already estimated in that proof. The Gram
conditioning, slope bound, and perturbation argument supply the
additional quantitative control needed here. Lean
checks the strengthened positive-integral estimate, slope growth, actual
Gram identities and conditioning, signed bilinear representation,
determinant comparison, and the complete quantitative assembly. The
original conditional minor-based route is retained in the source for
comparison, but is not used to discharge the final endpoint.
The elementary exponent-set lemmas and a Gram--Schmidt determinant lemma
are adapted, with attribution in the source, from OpenAI's September 2026
\emph{PiExponent} and \emph{Catalan} developments~\cite{OpenAI}. We do not claim those
general lemmas as new.

Throughout, $r\in\Q$ is fixed, $n\geq1$, and
\[
 \theta=\zeta(3)-r\zeta(2),\qquad h=3n,\qquad
 D_m(t)=\prod_{j=1}^m(t+j),\qquad
 R_n(t)=\frac{D_n(t)^4}{D_{5n}(t)}.
\]
Put
\[
 S_n=\frac{(5n)!}{(n!)^4},\quad
 F_n=\prod_{j=0}^{h-1}(j!)^2,\quad
 t(y)=\tfrac12+\ii y,
\]
\[
 w_r(y)=\frac{\pi}{2}\sech^2(\pi y)
        \bigl(2r-2\pi\ii\tanh(\pi y)\bigr),\qquad
 \phi_n(y)=t(y)R_n(t(y))w_r(y).
\]
The upstream rational matrices $A_{r,n},B_n$ satisfy
\begin{equation}\label{eq:entries}
 (A_{r,n}+\theta B_n)_{jk}
     =\int_\R t(y)^{j+k}\phi_n(y)\dd y,
       \qquad 0\leq j,k<h.
\end{equation}
Define
\[
 Q_n(X)=\det(A_{r,n}+XB_n),\qquad
 \widetilde Q_n(X)=\frac{S_n^h}{F_n}Q_n(X).
\]
There are $P_n\in\Z[X]$ and $d_n>0$ with
\begin{equation}\label{eq:arith}
 \deg P_n\leq3n,\qquad P_n=d_n\widetilde Q_n,
 \qquad \log d_n\leq(283/50+\varepsilon)n^2
\end{equation}
for all sufficiently large $n$ with $Q_n\ne0$, for every $\varepsilon>0$.
Here $d_n$ is the source's \texttt{dtilde}, not its \texttt{d}.

There is a finite exceptional prime set $E$, independent of $n$, such that
for every prime $p\geq7$ outside $E$ with $p\nmid\operatorname{den}(r)$,
\begin{equation}\label{eq:edge}
 P_{p-1}(X)\pmod p\quad\hbox{is a nonzero constant}.
\end{equation}
These inputs are respectively \texttt{arith\_node} and
\texttt{prime\_edge\_node} in \texttt{Zeta32/Final.lean}.

\section{Retaining the positive Heine integral}

Let
\[
 f_j(y)=\frac{t(y)^j}{j!},\qquad
 \rho_n(y)=S_n|\phi_n(y)|,\qquad
 M_n=\left(\int_\R \overline{f_j(y)}f_k(y)\rho_n(y)\dd y\right)_{j,k<h}.
\]
All polynomial moments exist: the rational factor has no real-line poles
on $t=1/2+\ii y$, has at most polynomial growth, and the hyperbolic kernel
decays exponentially. The density is positive except possibly at $y=0$.
Thus $M_n$ is a positive definite Hermitian matrix.

Andr\'eief's identity and the Vandermonde determinant give
\begin{equation}\label{eq:positive-integral}
 \det M_n=\frac{S_n^h}{F_nh!}
 \int_{\R^h}\prod_{\ell=1}^h|\phi_n(y_\ell)|
       \prod_{\ell<m}(y_\ell-y_m)^2\dd y.
\end{equation}
The factorial denominator is exactly $F_n$, since it occurs once in
each Vandermonde factor.

\begin{lemma}[Positive energy estimate]\label{lem:energy}
For every $\varepsilon>0$, eventually in $n$,
\[
                         \det M_n\leq e^{(-6+\varepsilon)n^2}.
\]
\end{lemma}
\begin{proof}
This is a strengthening of the \emph{conclusion} of the upstream analytic
node, not a consequence of its signed-determinant conclusion alone.
Here is the precise point at which the stronger quantity can be retained.
Write $I_n$ for the nonnegative integral in
\eqref{eq:positive-integral}, without its prefactor.
The proof of \texttt{heine\_scaled} in
\texttt{Zeta32/Analytic/Energy/Scaling.lean} establishes, for every integrable
majorant $D$ of its scaled configuration integrand,
\[
 \frac{S_n^h}{F_n}I_n
 \leq \exp\!\left(9(3/2-\log3)n^2+C_1n\log(n+1)\right)
       \int_{\R^h}D(x)\dd x.
\]
Specifically, its estimates \texttt{hscale}, \texttt{hpt}, \texttt{hint},
and the final coefficient estimate prove this inequality without using
the signed determinant. The factor $1/h!\leq1$ may be discarded.

The proof of \texttt{energy\_bound\_exp} in
\texttt{Zeta32/Analytic/Energy/Assembly.lean} then constructs a nonnegative
integrable $D$ with
\[
 \int D\leq
 \exp\!\left(9(\ell-I)n^2+C_2n(\log(n+1)+1)\right)K^{3n},
\]
where $K\geq1$ and the verified energy certificate gives
\[
                 9(3/2-\log3)+9(\ell-I)\leq-6.
\]
The error is $O_r(n\log(n+1))=o(n^2)$. Combining these estimates with
\eqref{eq:positive-integral} proves the assertion. The stronger integral
estimate has now been extracted formally in our \texttt{Zeta32Extension}
namespace as
\begin{center}\texttt{positiveHeine\_eventually\_le}.\end{center}
Its axiom closure contains only Lean's three standard axioms. Identification
with the actual Gram determinant is now checked by
\texttt{positiveGram\_det}.
\end{proof}

\section{An exponential lower bound for the Gram matrix}

\begin{lemma}[Weight on a distant interval]\label{lem:weight}
For $1\leq x\leq2$ and every integer $n\geq1$,
\[
                         n\rho_n(nx)\geq e^{-22n}.
\]
\end{lemma}
\begin{proof}
For $y\in[n,2n]$, every numerator factor of $R_n(t(y))$ has modulus
at least $n$, and every denominator factor has modulus at most $8n$.
The elementary factorial inequalities $m!\geq(m/e)^m$ and $n!\leq n^n$
therefore give
\[
 S_n|R_n(t(y))|\geq
       \left(\frac5{8e}\right)^{5n}\geq e^{-8n}.
\]
For the last inequality use $\log(8/5)\leq3/5$.
Also $|t(y)|\geq n$, $\tanh(\pi y)\geq1/2$, and
\[
 |w_r(y)|\geq\frac{\pi^2}{2}e^{-2\pi y}
             \geq e^{-4\pi n}.
\]
Indeed the imaginary part alone bounds the modulus of the second factor,
and $\cosh u\leq e^u$ for $u\geq0$. Since $8+4\pi<22$ and $n^2\geq1$,
the claimed bound follows, uniformly in $r$.
\end{proof}

\begin{lemma}[Polynomial coefficient recovery]\label{lem:poly}
For $h=3n$ and $v\in\C^h$, let
\[
 p(s)=\sum_{j=0}^{h-1}\frac{v_j}{j!}
                  \left(\tfrac12+\frac{\ii n(3+s)}2\right)^j.
\]
Then
\[
          \sum_{j<h}|v_j|^2\leq64h^5\,111132^h
                    \int_{-1}^1|p(s)|^2\dd s.
\]
\end{lemma}
\begin{proof}
Write $M=\max_{[-1,1]}|p|$. A real polynomial $u$ of degree at most $h$
and supremum at most $M$ on $[-1,1]$ has every coefficient bounded by
$7^hM$. Here is a coefficient proof of this standard Chebyshev estimate.
The extremal derivative inequality at $1$ bounds the absolute Taylor
coefficients of $u$ by those of $MT_h$, which are nonnegative. Translating
back from $1$ to $0$ costs at most $2^j$ for the $j$th Taylor coefficient.
Their sum is $MT_h(3)\leq M7^h$, the final inequality following from the
Chebyshev recurrence. Apply this to the real and imaginary coefficient
polynomials of $p$ to obtain $|[s^j]p|\leq2\,7^hM$.

For $s,t\in[-1,1]$, $|s^j-t^j|\leq j|s-t|$. Consequently
\[
 |p(s)-p(t)|\leq2h^2\,7^hM|s-t|.
\]
A one-sided interval of length $1/(4h^2 7^h)$ at a point where $|p|=M$
lies inside $[-1,1]$ and has $|p|\geq M/2$. Therefore
\[
 M^2\leq16h^2\,7^h\int_{-1}^1|p(s)|^2\dd s.
\]
The zero case is immediate. Define
\[
 q(z)=\sum_{j<h}\frac{v_jn^j}{j!}z^j
     =p(-3+\ii/n-2\ii z).
\]
The affine constant has modulus at most $4$ and the slope has modulus $2$.
Every coefficient of its $j$th power has modulus at most $6^j$; hence
$|[z^j]q|\leq2h42^hM$. Since $j!/n^j\leq3^h$ for $j<h=3n$,
we get $|v_j|\leq2h126^hM$. Square, sum, and use the preceding integral
estimate. The constant is $126^2\cdot7=111132$.
This proof, including the maximum-to-integral estimate, is checked in
\texttt{PolynomialConditioning.lean}.
\end{proof}

\begin{proposition}\label{prop:eigen}
For all sufficiently large $n$, uniformly in $r$, the smallest eigenvalue of $M_n$ satisfies
\[
                         \lambda_{\min}(M_n)\geq e^{-60n}.
\]
\end{proposition}
\begin{proof}
Restrict the Gram quadratic form to $y\in[n,2n]$, substitute
$y=n(3+s)/2$, and apply Lemmas~\ref{lem:weight} and \ref{lem:poly}:
\[
 v^*M_nv\geq\frac{e^{-22n}}{128h^5\,111132^h}\|v\|^2.
\]
Now $111132<e^{12}$ and $\log n=o(n)$ imply
$128(3n)^5\,111132^{3n}\leq e^{38n}$ eventually. This proves the result.
The full weighted-integral lower bound, the exact Gram quadratic-form
identity, and the resulting matrix-form inequality are checked in
\texttt{WeightedQuadratic.lean} and \texttt{GramQuadratic.lean}.
\end{proof}

\section{Slope growth and stability}

Normalize the entire pencil entrywise by
\[
 H_n(X)_{jk}=\frac{S_n}{j!k!}(A_{r,n}+XB_n)_{jk},\qquad
 T_n=H_n'(X).
\]
Then $\det H_n(X)=\widetilde Q_n(X)$. Notice that this determinant identity
holds for all $X$, not only at $\theta$.

\begin{lemma}[Normalized slope]\label{lem:slope}
For $0\leq j,k<3n$,
\[
                             |(T_n)_{jk}|\leq e^{50n}.
\]
\end{lemma}
\begin{proof}
An exact binomial identity gives, for $1\leq a\leq5n$,
\[
 S_n\operatorname{Res}_{t=-a}R_n(t)
  =(-1)^{a-1}5n\binom{5n-1}{a-1}\binom{a-1}{n}^{\!4}.
\]
This includes $a\leq n$, where both sides vanish. Since
$\binom uv\leq2^u$, the absolute value is at most $5n\,2^{25n}$.
The definition of $B_n$, followed by $(5n)^j/j!\leq e^{5n}$, gives
\[
 |(T_n)_{jk}|\leq250n^3\,2^{25n}
           \frac{(5n)^j}{j!}\frac{(5n)^k}{k!}
       \leq250n^3\,2^{25n}e^{10n}\leq e^{46n}\leq e^{50n}.
\]
For the penultimate step use $250\leq2^8\leq e^{8n}$, $n^3\leq e^{3n}$,
and $2^{25n}\leq e^{25n}$. This proof is now formalized in
\texttt{SlopeGrowth.lean}, with only the three standard axioms.
\end{proof}

There is also an independent contour derivation, useful for checking the
normalizations. The definition of $B_n$ is the sum of $2a$ times the residues of
$t^{j+k}R_n(t)$ at the poles $t=-a$, $1\leq a\leq5n$. Therefore, with
Laurent coefficients taken at infinity,
\[
 (B_n)_{jk}=-2[t^{-j-k-2}]R_n(t).
\]
This follows by applying the residue identity to $t^{j+k+1}R_n(t)$;
its polynomial part contributes no coefficient of $t^{-1}$.
On $|t|=6n$ the numerator factors have modulus at most $7n$, and the
denominator factors have modulus at least $n$. Cauchy's Laurent estimate
thus gives
\[
 |(B_n)_{jk}|\leq 2(6n)^{j+k+2}7^{4n}n^{-n}.
\]
Also $S_n\leq5^{5n}e^{4n}n^n$, and $(6n)^j/j!\leq e^{6n}$.
Hence
\[
 |(T_n)_{jk}|\leq72n^2
       e^{(12+5\log5+4+4\log7)n}\leq e^{50n}.
\]
For example $\log5,\log7<2$ and $72n^2\leq e^{6n}$ already give
the stronger intermediate bound $e^{40n}$.

\begin{lemma}[Bilinear Gram comparison]\label{lem:comparison}
Let $M$ be a positive definite Hermitian $h\times h$ matrix and let
$B$ be a complex $h\times h$ matrix. If $C\geq0$ and
\[
       |u^{\mathsf T}Bv|\leq C\sqrt{u^*Mu}\sqrt{v^*Mv}
                   \qquad(u,v\in\C^h),
\]
then $|\det B|\leq\det M\,C^h$.
\end{lemma}
\begin{proof}
Choose an invertible matrix $W$ with $W^*MW=I$, for example the inverse
positive square root. Set $L=W^{\mathsf T}BW$. The hypothesis implies
$|u^{\mathsf T}Lv|\leq C\|u\|_2\|v\|_2$. Taking
$u=\overline{Lv}$ gives $\|Lv\|_2^2\leq C\|Lv\|_2\|v\|_2$;
therefore $\|L\|_{2\to2}\leq C$, including the case $Lv=0$.
Hadamard's determinant inequality yields $|\det L|\leq C^h$.
Taking determinants of $W^*MW=I$ gives
$|\det W|^2\det M=1$, whereas
$|\det L|=|\det W|^2|\det B|$. These identities give the result.
The transpose in $W^{\mathsf T}BW$ is intentional: the original form
is bilinear, not sesquilinear. This lemma is formally checked as
\texttt{det\_le\_of\_bilinear\_gram}.
\end{proof}

For $F_v(y)=\sum_j v_jf_j(y)$, the exact signed representation is
\[
 u^{\mathsf T}H_n(\theta)v
       =\int_\R F_u(y)F_v(y)S_n\phi_n(y)\dd y.
\]
Since $|S_n\phi_n|=\rho_n$, the integral triangle inequality and
Cauchy--Schwarz give
\[
 |u^{\mathsf T}H_n(\theta)v|
       \leq\sqrt{u^*M_nu}\sqrt{v^*M_nv}.
\]
All the needed square integrals are finite by the moment estimates already
noted. Both this representation and this inequality are checked in
\texttt{SignedBilinear.lean}.

If a perturbation $E$ has entries bounded in modulus by $K\geq0$, then
finite Cauchy--Schwarz and Proposition~\ref{prop:eigen} imply
\[
 |u^{\mathsf T}Ev|\leq K\Bigl(\sum_j|u_j|\Bigr)
                          \Bigl(\sum_k|v_k|\Bigr)
       \leq hK e^{60n}\sqrt{u^*M_nu}\sqrt{v^*M_nv}.
\]
Apply Lemma~\ref{lem:comparison} to $B=H_n(\theta)+E$:
\begin{equation}\label{eq:perturb}
       |\det(H_n(\theta)+E)|\leq\det M_n(1+hKe^{60n})^h.
\end{equation}
This direct proof does not require any estimate for complementary minors.

\begin{proposition}[Local decay]\label{prop:local}
For fixed $r$, all sufficiently large $n$ with $Q_n\ne0$ satisfy
\[
 |x-\theta|\leq e^{-120n}
       \quad\Longrightarrow\quad |P_n(x)|\leq e^{-3n^2/10}.
\]
\end{proposition}
\begin{proof}
Take $\varepsilon=1/100$ in \eqref{eq:arith} and
Lemma~\ref{lem:energy}. In \eqref{eq:perturb} use
$E=H_n(x)-H_n(\theta)$ and $K=|x-\theta|e^{50n}$. Then
\[
 |\widetilde Q_n(x)|\leq
      e^{-5.99n^2}(1+3n e^{-10n})^{3n}
      \leq e^{-5.99n^2+3n},
\]
since $3n\leq e^{10n}$ and $2^{3n}\leq e^{3n}$. Multiplication by
$d_n\leq e^{5.67n^2}$ gives $|P_n(x)|\leq e^{-0.32n^2+3n}$.
For $n\geq150$, this is at most $e^{-0.30n^2}$.
\end{proof}

\section{Selecting a prime and extracting an exponent}

An elementary Chebyshev estimate gives, eventually in real $u$,
\[
       \sum_{u<p\leq3u}\log p>u/4.
\]
Our Catalan development proves this as
\texttt{eventually\_primeLogInterval\_gt\_quarter}; no prime number
theorem is needed at this step. Put $u=20\log q$. If all primes in this
window divided $q$, their product would divide $q$ and the displayed
sum would be at most $\log q$, contradicting its lower bound
$5\log q$. Any fixed finite prime exclusions can be placed below the
window by increasing the threshold for $q$.

\begin{theorem}[Quantitative extension]\label{thm:bound}
For every fixed $r\in\Q$, there exists $q_0(r)$ such that every reduced
rational number $a/q$, with $a\in\Z$ and $q\in\mathbb N$, satisfies
\[
 \left|\zeta(3)-r\zeta(2)-\frac aq\right|>q^{-10000}
       \qquad(q\geq q_0(r)).
\]
In particular,
$2\leq\mu(\zeta(3)-r\zeta(2))\leq10000$.
\end{theorem}
\begin{proof}
Choose a prime $p\nmid q$ with
$20\log q<p\leq60\log q$, above all fixed cutoffs for
\eqref{eq:edge} and Proposition~\ref{prop:local}; put $n=p-1$.
Eventually,
\[
                     19\log q<n\leq60\log q.
\]
In particular $p\nmid\operatorname{den}(r)$.
If $|\theta-a/q|\leq q^{-10000}$, then
\[
 q^{-10000}\leq e^{-120n},
 \qquad
 |P_n(a/q)|\leq e^{-3n^2/10}<q^{-3n}.
\]
For the strict last inequality it suffices that $n>10\log q$.
On the other hand, \eqref{eq:edge} implies $P_n(a/q)\ne0$, since $q$ is
invertible modulo $p$. As $\deg P_n\leq3n$, the number $q^{3n}P_n(a/q)$
is a nonzero integer. Thus $|P_n(a/q)|\geq q^{-3n}$, a contradiction.
The lower bound $2$ follows from the upstream irrationality theorem and
Dirichlet's approximation theorem.
\end{proof}

\section{Formal verification and scope}

This argument proves a finite exponent with explicit numerical constant,
using the cited upstream construction. No uniformity in the threshold
$q_0(r)$ as the height of $r$ grows is asserted. Consequently
it is not a simultaneous linear-independence measure for arbitrary varying
coefficients of $1,\zeta(2),\zeta(3)$. The case $r=0$ already has classical,
much better finite bounds; the intended new targets have $r\ne0$.

All parts of the direct proof compile in Lean. The final theorem names are
\begin{center}
\texttt{Zeta32Extension.denominator\_bound},\\
\texttt{Zeta32Extension.irrationalityExponent\_bounds}.
\end{center}
Their only parameter is $r:\Q$; no analytic estimate is assumed. An
explicit audit checks all 123 extension theorems and permits only
\texttt{propext}, \texttt{Classical.choice}, and \texttt{Quot.sound}.
A fresh-kernel replay passed for 21 selected roots and their full closure
of 82674 declarations, including both unconditional endpoints.
The originally considered theorem
\begin{center}\texttt{denominator\_bound\_of\_uniformMinors}\end{center}
remains conditional in the
source, but is not needed to discharge the final bound. Merely auditing
that conditional statement would not have proved the present result.

Reproduction commands and pinned dependency versions are provided in
\texttt{anc/README.txt} in the source archive. The replay uses Lean's own
kernel. The argument does not compute a numerical value of $q_0(r)$.

Dauguet--Zudilin (2014)~\cite{DZ},
\href{https://arxiv.org/abs/1401.5322}{\emph{On simultaneous diophantine
approximations to $\zeta(2)$ and $\zeta(3)$}}, bounds certain restricted
linear forms, not an unrestricted finite exponent for each such difference.
The recent qualitative result of Tang is an essential acknowledged input.

\paragraph{Acknowledgement.}
The mathematical development, Lean formalization, and preparation of this
manuscript were carried out with substantial assistance from OpenAI Codex.

\begin{thebibliography}{9}
\raggedright
\bibitem{Tang}
Q. Tang, \emph{zeta2-zeta3-linear-independence}, Lean development,
September 2026. Pinned revision specified in Section~1.
\url{https://github.com/dtq1997/zeta2-zeta3-linear-independence}.

\bibitem{DZ}
S. Dauguet and W. Zudilin,
\emph{On simultaneous diophantine approximations to $\zeta(2)$ and
$\zeta(3)$}, 2014. \url{https://arxiv.org/abs/1401.5322}.

\bibitem{OpenAI}
OpenAI, \emph{math}: formal developments accompanying
\emph{Catalan's constant is irrational} and
\emph{The irrationality exponent of pi is 2}, September 2026.
\url{https://github.com/openai/math}.
Only the elementary approximation-exponent and Hilbert determinant lemmas
identified in the source attribution are adapted here; neither endpoint
theorem is used as an assumption.
\end{thebibliography}
\end{document}
